% =====================================================================
%  18.2  Range Variants: Differences, Pairs, Compression   sec:c14-range
%  Exports: def:c14-diff  thm:c14-twobit  prop:c14-compress
%           prop:c14-prefixmin  lst:c14-diff-build  lst:c14-diff-range
%           lst:c14-twobit  lst:c14-compress  lst:c14-min  lst:c14-1651-soln
%           prob:c14-rangeupd
%  Shared numbers (all hand-reverified, ASCII only):
%    D = (3,-1,2,1,-4,0,4,-2); diff-BIT Td = BIT over D = (3,2,2,5,-4,-4,4,3);
%    add +1 to [2,5]: Td -> (3,3,2,6,-4,-5,4,3); reads: p4=6, p2=3, p7=5, p8=3.
%    Two-BIT (T1 over D, T2 over j*D_j), prefix(x)=(x+1)Q1(x)-Q2(x):
%    original T1=(3,2,2,5,-4,-4,4,3), T2=(3,1,6,11,-20,-20,28,3);
%    prefix(4)=5*5-11=14, prefix(6)=7*1+9=16, prefix(3)=4*4-7=9;
%    after +1 on [2,5]: T1=(3,3,2,6,-4,-5,4,3), T2=(3,3,6,13,-20,-26,28,-1);
%    sum[1,4]=5*6-13=17, sum[1,6]=7*1+13=20, sum[4,6]=20-11=9.
%  CSES 1651 sample (live-verified): 8 3 / 3 2 4 5 1 1 5 3 /
%    2 4 / 1 2 5 1 / 2 4  ->  outputs 5 then 6 (x[4]=5; after +1 on [2,5], 6).
%  Two-BIT convention used EVERYWHERE here: T1 over D, T2 over j*D_j, and
%    prefix(x) = (x+1)Q1(x) - Q2(x)   (thm:c14-twobit, lst:c14-twobit).
%    Verified on a=(3,2,4,5,1,1,5,3): prefix(4)=5*5-11=14=sum a[1..4];
%    after +1 on [2,5]: T1=(3,3,2,6,-4,-5,4,3), T2=(3,3,6,13,-20,-26,28,-1),
%    sum[1,4]=5*6-13=17; sum[4,6]=20-11=9 (see the side-by-side tracebox).
% =====================================================================
\section{Range Variants: Adding to Intervals, Summing Intervals}\label{sec:c14-range}

\begin{previouslybox}[18.1 stores sums of \emph{data}; now store sums of \emph{changes}]
The plain BIT answers a prefix of \(A\) and mutates one entry of \(A\).
This section swaps which side gets the expensive operation: encode the
\emph{difference array} \(D_i=A_i-A_{i-1}\) in the BIT, and a whole-interval
addition collapses to two point updates -- at the price that a read becomes
a prefix walk.  Push the algebra one step further (a \emph{paired} BIT) and
even ``range add \emph{and} range sum'' become logarithmic.  Two supporting
tools travel with them: coordinate compression (when indices are not small
integers) and the prefix-min variant (when the accumulator is \(\min\), not
\(\sum\)).
\end{previouslybox}

\begin{topicintro}
\textbf{Prerequisites.} \secref{c14-core} in full -- walks, blocks,
\(\coderef{c14-build}\).  The telescoping identity
\(A_i=\sum_{j\le i}(A_j-A_{j-1})\) (Ch.~5 used it for edge-span counts).

\textbf{After this section you can}
\begin{itemize}[nosep]
  \item build the difference BIT and answer range-add/point-read inputs
        (CSES~1651); \hfill\lv{3}--\lv{4}
  \item derive and implement the two-BIT range-add/range-sum identity from
        the contribution count, not from memory; \hfill\lv{4}--\lv{5}
  \item coordinate-compress arbitrary values into BIT indices without
        breaking equal-range semantics; \hfill\lv{4}
  \item state exactly when a prefix-min BIT is legal (append-only decreases)
        and when it silently lies. \hfill\lv{4}
\end{itemize}
\end{topicintro}

\begin{recognitionbox}[Pattern card -- ``intervals move, positions/segments are asked'']
\reckey{``increase each value in \([a,b]\) by \(u\)'' + ``value at position
\(k\)'' (CSES~1651's exact words); ``add to a range, then ask a range sum''
(two BITs); values up to \(10^9\) but only \(\le2\cdot10^5\) of them appear
(compression); ``\(\mathit{dp}[i]=\min_{j<i,\;\text{cond}}\mathop(\ldots)\)
with \(j\) queried by a prefix of \(\min\) values'' (prefix-min BIT).}
\recuse{
  \item range add, point read -- difference BIT, this section (row~3 of the
        opener's cost chart);
  \item range add \emph{and} range sum -- two paired BITs
        (\thmref{c14-twobit}); the segment tree with lazy (\chref{c15}) also
        does it, at 4x the code;
  \item index domain huge, order the only thing that matters -- compress
        first (\propref{c14-compress});
  \item accumulator \(\min\) with monotone (decreasing-only) writes --
        prefix-min BIT (\propref{c14-prefixmin}).
}
\recnotuse{
  \item point \emph{sets} and range sums -- the plain BIT of
        \secref{c14-core} (differences would force you to re-point two
        slots per set and buy nothing);
  \item arbitrary overwrites of a \(\min\) table -- the monotonicity
        hypothesis of \propref{c14-prefixmin} fails; go to \chref{c13};
  \item range \emph{assign} (set all of \([a,b]\) to \(v\)) -- neither
        difference nor pair survives non-invertible assignment; lazy Ch.~20.
}
One-line test: \emph{if every mutation touches a whole interval, something
telescoping is being stored -- put \emph{that} in the BIT.}
\end{recognitionbox}

% ---------------------------------------------------------------
\subsection{The Difference BIT: Range Add, Point Read}
% ---------------------------------------------------------------

Let \(D_1=A_1\), \(D_i=A_i-A_{i-1}\).  Then \(A_i=\sum_{j\le i}D_j\): a
\emph{point read of \(A\) is a prefix query of \(D\)} -- exactly the
operation a BIT is born for.  The magic is the other direction:

\begin{definition}{Difference encoding of interval adds}{c14-diff}
Adding \(\delta\) to every \(A_p\) with \(p\in[a,b]\) changes only two
differences: \(D_a\mathrel{+}=\delta\) (the step up into the interval) and
\(D_{b+1}\mathrel{-}=\delta\) (the step back out).  All other \(D_j\) are
untouched because inside the interval both neighbours moved by \(\delta\).
\end{definition}

So the whole workload ``range add / point read'' becomes: BIT over \(D\);
\(\texttt{add}(a,\delta),\ \texttt{add}(b{+}1,-\delta)\); read
\(\code{query}(k)\).  \(\bigO{\log n}\) per op, three lines of code.

\begin{lstlisting}[caption={the difference BIT: build from A by differencing (O(n)), range add as two point updates}, label={lst:c14-diff-build}]
// t[] is a BIT over D; every "A[i]" below is 1-based
void build_diff(){                       // O(n): t[i] = D[i] first...
    for (int i = 1; i <= n; i++)
        t[i] = A[i] - (i > 1 ? A[i-1] : 0);
    for (int i = 1; i <= n; i++){        // ...then the same push as
        int j = i + lowbit(i);           // coderef:c14-build
        if (j <= n) t[j] += t[i];
    }
}
void range_add(int a, int b, long long d){
    update(a, d);                        // t is the point-update walk of
    if (b + 1 <= n) update(b+1, -d);     // coderef:c14-update on D
}
long long point_read(int k){ return query(k); }   // prefix of D = A[k]
\end{lstlisting}

\begin{tracebox}[Difference cascade on the CSES 1651 sample array]
\(a=(3,2,4,5,1,1,5,3)\) gives the raw difference array
\(D=(3,-1,2,1,-4,0,4,-2)\).  The BIT is built on top of \(D\) by
\(\coderef{c14-diff-build}\): start \(t=D\) and push \(t[i]\to t[i{+}\text{lowbit}(i)\)
(\(\coderef{c14-build}\)), giving the \emph{block sums}
\[
  Td=(\,3,2,2,5,-4,-4,4,3\,),
\]
cell by cell: \(Td[1]=3\), \(Td[2]=D_1{+}D_2=2\), \(Td[3]=D_3=2\),
\(Td[4]=\sum_{1..4}D=5\), \(Td[5]=D_5=-4\), \(Td[6]=D_5{+}D_6=-4\),
\(Td[7]=D_7=4\), \(Td[8]=\sum_{1..8}D=3\).  This is \(Td\), \emph{not} the
raw \(D\) above -- conflating them is Pitfall~1 below.

Now \(\texttt{range\_add}(2,5,+1)\) fires \(\mathit{update}(2,+1)\) (walk
\(2{\to}4{\to}8\): \(Td[2]\!{\to}\!3,\ Td[4]\!{\to}\!6,\ Td[8]\!{\to}\!4\)) and
\(\mathit{update}(6,-1)\) (walk \(6{\to}8\): \(Td[6]\!{\to}\!-5,\ Td[8]\!{\to}\!3\)),
so the live tree is \(Td=(3,3,2,6,-4,-5,4,3)\).  Point reads are prefix walks:
\(\text{read}(4)=Td[4]=6=a_4{+}1\) \checkmark,\quad
\(\text{read}(2)=Td[2]=3=a_2{+}1\) \checkmark,\quad
\(\text{read}(7)=Td[7]{+}Td[6]{+}Td[4]=4-5+6=5=a_7\) \checkmark\ (position~7
lies outside \([2,5]\)); \(\text{read}(8)=Td[8]=3=a_8\) is unchanged because
the two difference writes contributed \(+1-1=0\) to \(Td[8]\) -- the encoding
cancels exactly as \defref{c14-diff} promises.
\end{tracebox}

% ---------------------------------------------------------------
\subsection{Two BITs: Range Add \emph{and} Range Sum}
% ---------------------------------------------------------------

Difference reads a point; but a \emph{sum} of \(A\)'s prefix in terms of
\(D\) weights each difference by how many entries it still affects:
\[
  \sum_{i=1}^{x} A_i
  = \sum_{i=1}^{x}\sum_{j\le i} D_j
  = \sum_{j\le x} D_j\,(x-j+1)
  = (x+1)\underbrace{\sum_{j\le x} D_j}_{S_1(x)}
    \;-\; \underbrace{\sum_{j\le x} j\,D_j}_{S_2(x)} .
\]
Two plain BITs -- one over \(D_j\), one over \(j\,D_j\) -- answer range
add/range sum in \(4\) updates and \(4\) walks per op.

\begin{theorem}{The two-BIT identity}{c14-twobit}
Maintain BITs \(T_1\) over \(D\) and \(T_2\) over \(j\cdot D_j\).  Then
\(\displaystyle\mathit{prefix}(x)=(x+1)\,\code{query}_1(x)-\code{query}_2(x)\)
and a range add \(\delta\) on \([a,b]\) is
\[
\text{at }a:\ \mathit{add_1}(a,\delta),\ \mathit{add_2}(a,a\delta);
\qquad
\text{at }b{+}1:\ \mathit{add_1}(b{+}1,-\delta),\ \mathit{add_2}(b{+}1,-(b{+}1)\delta).
\]
Every operation is \(\bigO{\log n}\).
\end{theorem}
\begin{proofbox}
The prefix formula is the telescoping rearrangement above -- each \(D_j\)
lifts \(A_j,\dots,A_x\), i.e.\ \(x-j+1\) entries, so its contribution to
\(\sum_{i\le x}A_i\) is \((x+1)D_j-jD_j\).  A range add changes exactly
\(D_a{+}=\delta\), \(D_{b+1}{-}=\delta\) (\defref{c14-diff}); each change
must be entered in \emph{both} trees with the respective weight
(\(\delta\) vs.\ \(j\delta\)), giving the four updates.  \code{query}
correctness is \thmref{c14-query} applied to each tree.
\end{proofbox}

\begin{lstlisting}[caption={range add + range sum with two BITs (worked example, not a linked problem: CSES has no such task)}, label={lst:c14-twobit}]
// T1 over D, T2 over j*D[j]; both use coderef:c14-update/coderef:c14-query
void ra(int i, long long d){                 // one difference point change
    for (int k = i; k <= n; k += lowbit(k)){
        t1[k] += d;                          // d contribution
        t2[k] += d * i;                      // (j * d) contribution, j = i
    }
}
void range_add(int a, int b, long long d){
    ra(a, d);  ra(b + 1, -d);
}
long long bitq(long long* t, int i){
    long long s = 0;
    for (; i > 0; i -= lowbit(i)) s += t[i];
    return s;
}
long long prefix(int x){ return (x + 1) * bitq(t1, x) - bitq(t2, x); }
long long range_sum(int a, int b){ return prefix(b) - prefix(a - 1); }
\end{lstlisting}

\begin{tracebox}[Both BITs side by side: start a=(3,2,4,5,1,1,5,3), add +1 to [2,5], then sum [1,4] and [4,6]]
Before any update, \(D=(3,-1,2,1,-4,0,4,-2)\).  \(T_1\) is the BIT of block
sums over \(D\) -- exactly the \(Td\) row of the box above,
\(T_1=(3,2,2,5,-4,-4,4,3)\).  \(T_2\) is the BIT over \(jD_j=(3,-2,6,4,-20,0,28,-16)\),
which \(\coderef{c14-build}\) turns into \(T_2=(3,1,6,11,-20,-20,28,3)\).
Sanity check on the untouched array: \(\mathit{sum}[1,4]=\mathit{prefix}(4)=5{\cdot}Q_1(4)-Q_2(4)=5{\cdot}5-11=14=3{+}2{+}4{+}5\) \checkmark.

Apply \(\texttt{range\_add}(2,5,+1)\): four point writes -- at \(j=2\) put
\((\delta,\,j\delta)=(+1,+2)\) into \((T_1,T_2)\), at \(j=6\) put \((-\delta,\,-j\delta)=(-1,-6)\).
Their update walks give \(T_1=(3,3,2,6,-4,-5,4,3)\) and \(T_2=(3,3,6,13,-20,-26,28,-1)\).
Now the range sums read off the live pair:
\[
\mathit{sum}[1,4]=\mathit{prefix}(4)=5{\cdot}T_1[4]-T_2[4]=5{\cdot}6-13=17=(3{+}3{+}5{+}6)\ \checkmark ,
\]
\[
\mathit{sum}[4,6]=\mathit{prefix}(6)-\mathit{prefix}(3)=(7{\cdot}1+13)-(4{\cdot}5-9)=20-11=9=(6{+}2{+}1)\ \checkmark .
\]
Both trees moved by the \emph{same} two difference events; the second simply
weights each by its index \(j\), which is what turns a point read into a range
sum.  This is \(\thmref{c14-twobit}\) running, no memorised constant required.
\end{tracebox}

% ---------------------------------------------------------------
\subsection{Coordinate Compression}
% ---------------------------------------------------------------

A BIT is an array indexed by \(1..n\), so it silently assumes the things you
store live at small \emph{integer} positions.  When the real indices are huge
(values up to \(10^9\)) or non-integral, only their \emph{relative order} can
matter to a prefix count -- and order is exactly what survives being replaced
by rank.  Compression is the one-time relabelling that makes that legal.

\begin{proposition}{Compression preserves the game}{c14-compress}
Let \(\mathit{vals}\) be the sorted, deduplicated set of all values that will
\emph{ever} be indexed (read, updated, or compared against).  Replacing each
value by its index in \(\mathit{vals}\) turns the value domain into
\([1..m]\), \(m\le\) (number of events), and preserves every comparison
\(v_1\le v_2\), hence every prefix count/sum the BIT answers.  Cost:
\(\bigO{M\log M}\) once.
\end{proposition}

\begin{lstlisting}[caption={standard compression: collect, sort, unique, lower\_bound}, label={lst:c14-compress}]
// Compress BEFORE any update/query: every value that will ever be an index
// must already sit in vals when the sort happens.  See thm:c14-twobit for a
// value-weighted tree when magnitudes (not just order) are needed.
vector<int> vals;                          // 1. COLLECT every value that
for (...) vals.push_back(x);               //    will ever be an index
sort(vals.begin(), vals.end());            // 2. SORT + 3. UNIQUE
vals.erase(unique(vals.begin(), vals.end()), vals.end());
int rank_(int x){                          // 4. RANK by lower_bound
    return (int)(lower_bound(vals.begin(), vals.end(), x) - vals.begin()) + 1;
}
// rank_ returns a 1-based index into vals -- the BIT's [1..m] domain.
// from here on, the BIT is sized m = vals.size(), never 1e9
\end{lstlisting}

% ---------------------------------------------------------------
\subsection{Prefix Min, and the Only Way It Is Legal}
% ---------------------------------------------------------------

\begin{proposition}{The append-only prefix-min BIT}{c14-prefixmin}
If the accumulator is \(\min\) and slot values only ever \emph{decrease} (the
standard DP pattern \(\mathit{dp}[i]=\min_{j\le i}\) over processed prefixes,
each \(j\) written once), then storing block minima and updating
\(t[k]=\min(t[k],v)\) along the \emph{update} walk answers
\(\min_{j\le x}\) along a \emph{superset-cover} walk: query visits the same
\(i\to i{-}\operatorname{lowbit}(i)\) blocks, which cover \([1..x]\) but now
may also be joined by over-covering ancestors -- for \(\min\) (idempotent)
extra coverage is harmless.  Arbitrary point increases are \textbf{not}
supported: the stale minimum remains buried in some slot with no rule for
which one to fix.
\end{proposition}

\begin{lstlisting}[caption={prefix-min BIT under monotone writes (query min over [1..x])}, label={lst:c14-min}]
void push(int i, long long v){             // position i may only decrease
    for (; i <= n; i += lowbit(i)) t[i] = min(t[i], v);
}
long long qmin(int x){                     // min over [1..x] (over-covers)
    long long m = INF;
    for (; x > 0; x -= lowbit(x)) m = min(m, t[x]);
    return m;
}
\end{lstlisting}

\begin{warningbox}[The prefix-min BIT answers min only while writes stay monotone]
``Cover a superset, harmless for \(\min\)'' is exactly as true for \(\max\);
what is \emph{not} true is mixing directions -- a \(\min\)-BIT whose writes
sometimes raise a value has no local repair rule (the raised entry may hide
in any ancestor).  When monotonicity breaks, you need recomputation over
blocks: that is a segment tree, \chref{c13}, and the honest cost is stating
so here instead of debugging a silent WA for a night.
\end{warningbox}

% ---------------------------------------------------------------
\subsection{BIT vs Segment Tree: the Boundary in One Table}
% ---------------------------------------------------------------

\begin{comparisonbox}[Where each range engine lives (and why this book teaches both)]
\begin{itemize}[nosep,leftmargin=1.4em]
  \item \textbf{BIT}: sums/counts (invertible) or monotone \(\min/\max\)
        writes; point or difference updates; \(\sim\)8 lines per flavour;
        trivially 2D-able (\(\secref{c14-offline}\), practice CF~1093G);
        cache-friendly half the memory of a segtree.
  \item \textbf{Segment tree} (Ch.~19, \chref{c13}): \emph{any} associative
        aggregate (max/min/gcd under arbitrary writes), non-invertible
        queries, custom combine functions; \(\sim\)40 lines; lazy variant
        (Ch.~20, \chref{c15}) for range assign/add with range asks.
  \item \textbf{Rule of thumb for contests}: if you can phrase every query
        as ``prefix of an invertible accumulator'' or a monotone DP write,
        the BIT is the lower-effort AC; the moment the statement asks a
        \emph{suffix min under increases}, \emph{range mode}, or
        \emph{first position where a non-additive predicate flips}, reach
        for \chref{c13}.
\end{itemize}
This is prose, per the chapter contract -- segtree code appears only there.
\end{comparisonbox}

% ---------------------------------------------------------------
\subsection{Problem Lab: CSES 1651 Range Update Queries}
% ---------------------------------------------------------------

\solvedproblem{c14-rangeupd}{Range Update Queries}{CSES 1651}
  {https://cses.fi/problemset/task/1651}{CSES L4}
  {difference BIT, range add as two point updates, point read as prefix}

\begin{problemstatement}[CSES 1651 Range Update Queries]
An array of \(n\) values; process \(q\) queries of two types:
\texttt{1 a b u} --- increase each value in \([a,b]\) by \(u\);
\texttt{2 k} --- print the value at position \(k\).

\textbf{Input.} \(n,q\le 2\cdot10^5\); \(|x_i|\le10^9\); \(|u|\le10^9\)
(negative updates allowed); \(1\le a\le b\le n\), \(1\le k\le n\).

\textbf{Output.} Each type-2 query on its own line.

\textbf{Example.}
\begin{verbatim}
Input:
8 3
3 2 4 5 1 1 5 3
2 4
1 2 5 1
2 4

Output:
5
6
\end{verbatim}
\end{problemstatement}

\begin{attemptgate}[Attempt it yourself first -- \lv{4}, give yourself 25 minutes]
The type-1 brute force loops \(a..b\): a single \([1,n]\) update is
\(\bigO{n}\), and \(q\) of them is \(\tle\).  The hint is in what type-1
\emph{does not} change: differences between adjacent values.
\peek{after adding \(+1\) to \([2,5]\), which adjacent pairs actually see a
change in \(A_i-A_{i-1}\)?}
\nudge{exactly two boundaries: at \(2\) (enters) and \(6\) (leaves) -- store
the difference array in the BIT of \secref{c14-core} and you already know
every remaining line of code.}
\fullhint{BIT over \(D\).  \texttt{1 a b u}: \code{update(a,u)},
\code{update(b+1,}-u\code{)} (guard \(b+1\le n\)).  \texttt{2 k}:
\code{query(k)} \(=\sum_{j\le k}D_j=A_k\).  Build: \(t[i]\!\leftarrow\!D_i\)
then the push loop of \(\coderef{c14-diff-build}\).  \code{long long}
everywhere: \(|u|\le10^9\), \(2\cdot10^5\) overlapping adds reach
\(2\cdot10^{14}\).}
\end{attemptgate}

\solstep{Understanding}
The problem never asks a range sum -- it asks \emph{single positions} after
\emph{whole-interval} moves.  That is the difference encoding's exact
signature (\defref{c14-diff}): the mutation is local in \(D\) (two slots),
and the read is global-in-prefix in \(D\) (one walk).  A plain BIT over
\(A\) would have it backwards: reads \(O(\log)\), writes \(O(n)\).

\solstep{Key Observations}
\begin{observation}
An interval add is a \emph{rank-one} change in the difference view: \(+u\) at
the left edge, \(-u\) at one past the right edge.  Everything inside shifts
together, so no difference between consecutive inside-positions changes.
\end{observation}
\begin{observation}
The initial build must difference \emph{first}, then BIT-build (\(\coderef{c14-diff-build}\));
feeding \(A\) into \(\coderef{c14-build}\) and then calling updates
``range-style'' answers \(A_k\) with a sum of raw values, not the value --
the Pitfall~1 trace above shows how plausible this looks.
\end{observation}
\begin{observation}
\(b=n\) is safe: \(\mathit{update}(n{+}1,\cdot)\) is out of the array's
domain \([1,n]\) and no read walk ever passes \(n\); guard with
\(b{+}1\le n\) (or size \(t\) to \(n{+}1\) and skip the guard).
\end{observation}

\solstep{Think It Through With Me}
\begin{thinkbox}
\deadend{Lazy segment tree (Ch.~20's tool)}
{Correct and \(\bigO{\log n}\) per op, but ~45 lines, \(4n\) memory, and a
propagation discipline to debug -- for ``range add, point read'' the lazy
machinery answers a question nobody asked (range aggregates).  Reserve it
for \chref{c15} proper.}
\deadend{Difference \emph{array} (no BIT)}
{The static array version supports each update in \(O(1)\) but each read in
\(\bigO{n}\) (re-prefix the whole array).  Interleaving is what kills it --
\(q\) reads after \(q\) updates is the same \(\bigO{nq}\).  The BIT is what
makes the prefix-of-differences incremental.}
\aha{Store in the tree what makes the \emph{update} cheap, if the \emph{query}
is an operation the tree already answers fast -- here differences turn
intervals into points, and prefix sums are the BIT's native tongue.}
\end{thinkbox}

\solstep{The Approach Ladder}
\begin{center}\renewcommand{\arraystretch}{1.25}
\begin{tabular}{@{}c p{0.235\linewidth} p{0.30\linewidth} l c@{}}
\toprule
\textbf{\#} & \textbf{Approach} & \textbf{Idea} & \textbf{Time} & \textbf{Verdict}\\
\midrule
1 & \lv{2}\ loop the interval & \(A_p+=u\) for \(p\in[a,b]\) & \(\bigO{nq}\) & \tle \\
2 & \lv{3}\ diff array only & two \(O(1)\) writes, re-prefix per read & \(\bigO{nq}\) & \tle \\
3 & \lv{4}\ \textbf{difference BIT} & two \(\bigO{\log}\) writes, \(\bigO{\log}\) read & \(\bigO{(n+q)\log n}\) & \ac \\
4 & \lv{4}\ lazy segtree & full power, mostly unused here & \(\bigO{(n+q)\log n}\) & overkill \\
\bottomrule
\end{tabular}\end{center}

\solstep{Solution}
Build \(D\) from the input, then the \(\bigO{n}\) BIT build
(\(\coderef{c14-diff-build}\)).  Type~1: two point updates at \(a\) and
\(b{+}1\).  Type~2: \code{query(k)}.  \code{long long} for the tree and the
answers.

\solstep{Proof of Correctness}
\begin{proofbox}
\textbf{Invariant.} After any sequence of processed queries, \(t\) is a BIT
over the current difference array \(D\) of the logical array \(A\).
\emph{Init:} \(\coderef{c14-diff-build}\) establishes it (\propref{c14-linear}
applied to \(D\)).  \emph{Type-1 step:} by \defref{c14-diff} the interval add
changes exactly \(D_a,D_{b+1}\) by \(+u,-u\); the two \code{update} walks
maintain the BIT for those point changes (\thmref{c14-query}).
\emph{Type-2 step:} \(A_k=\sum_{j\le k}D_j\) telescopes; \code{query(k)}
returns that prefix by \thmref{c14-query}'s disjoint-cover property.
\textbf{Bounds:} values \(|A_k|\le 10^9+2\cdot10^5\cdot10^9<3\cdot10^{14}\)
-- \code{long int long}; walks \(\le18\). \textbf{Termination:} each walk
strictly increases/decreases its index past the guards.
\end{proofbox}

\solstep{Implementation}
\begin{lstlisting}[caption={CSES 1651 full solution (difference BIT)}, label={lst:c14-1651-soln}]
#include <bits/stdc++.h>
using namespace std;
const int MX = 200005;
int n, q;
long long t[MX];

int lowbit(int i){ return i & -i; }
void update(int i, long long d){
    for (; i <= n; i += lowbit(i)) t[i] += d;
}
long long query(int i){
    long long s = 0;
    for (; i > 0; i -= lowbit(i)) s += t[i];
    return s;
}
int main(){
    scanf("%d%d", &n, &q);
    long long prev = 0, cur;
    for (int i = 1; i <= n; i++){          // build D in place
        scanf("%lld", &cur);
        t[i] = cur - prev; prev = cur;     // t = raw D so far
    }
    for (int i = 1; i <= n; i++){          // O(n) BIT push (build_diff)
        int j = i + lowbit(i);
        if (j <= n) t[j] += t[i];
    }
    while (q--){
        int ty, a, b; long long u;
        scanf("%d", &ty);
        if (ty == 1){
            scanf("%d%d%lld", &a, &b, &u);
            update(a, u);
            if (b + 1 <= n) update(b + 1, -u);
        } else { scanf("%d", &a); printf("%lld\n", query(a)); }
    }
}
\end{lstlisting}

\solstep{Trace on the Official Sample}
\begin{tracebox}[CSES 1651: read x[4], add +1 to [2,5], read x[4] again]
\(D=(3,-1,2,1,-4,0,4,-2)\); after the push phase the live tree is
\(Td=(3,2,2,5,-4,-4,4,3)\) (the row verified in the first trace box of this
section).
\[
\renewcommand{\arraystretch}{1.15}
\begin{array}{c|l|l|l}
\text{\#} & \text{query} & \text{walk} & \text{result}\\ \hline
1 & \texttt{2 4} & \text{query}(4)=Td[4]=5 & 5\\
2 & \texttt{1 2 5 1} & \text{update}(2,+1):\ Td[2]{\to}3,\ Td[4]{\to}6,\ Td[8]{\to}4 & \text{(add)}\\
  & \text{(cont.)} & \text{update}(6,-1):\ Td[6]{\to}-5,\ Td[8]{\to}3 & \\
3 & \texttt{2 4} & \text{query}(4)=Td[4]=6 & 6
\end{array}
\]
Printed: \texttt{5, 6} -- matches the live statement (each type-2 answer on
its own line).  Cross-check on \(\text{query}(8)=Td[8]=3=a_8\) after the
update: \(a_8=3\) was outside \([2,5]\), and \(Td[8]\) received \(+1-1=0\)
net -- the difference encoding cancelling itself exactly as
\defref{c14-diff} promises.
\end{tracebox}

\solstep{Complexity}
\begin{complexitybox}
\(\bigO{n}\) build, \(\bigO{\log n}\) per update/read; total
\(\lesssim 2\cdot10^5 + 3\cdot10^5\cdot18\approx5.6\cdot10^6\) ops.
Memory one \code{long long} array.  \textbf{AC}.
\end{complexitybox}

\solstep{Verification and Edge Cases}
\begin{center}\renewcommand{\arraystretch}{1.2}
\begin{tabular}{@{}p{0.20\linewidth} p{0.30\linewidth} p{0.38\linewidth}@{}}
\toprule
\textbf{Case} & \textbf{Why tricky} & \textbf{Code behaviour}\\
\midrule
\(b=n\) & \(D_{n+1}\) has no slot & guard \code{b+1<=n}; no read ever reaches \(n{+}1\) \\
negative \(u\) & tree entries go negative & \code{long long}, no logic change \\
overlapping adds & value drifts to \(2\cdot10^{14}\) & magnitudes bounded above; fits \\
first element & \(D_1=A_1\) (no \(A_0\)) & \code{prev=0} initialisation handles it \\
read after many adds & needs \emph{current} value & prefix walk reads live tree \\
\bottomrule
\end{tabular}\end{center}

\begin{hardeningbox}
\begin{itemize}[nosep,leftmargin=1.4em]
  \item \textbf{Do not mix flavours:} a tree built over \(D\) answers point
        reads; feeding it ``set \(A_k=u\)'' queries (as in CSES~1648) is
        \emph{wrong} -- a set at \(k\) is two delta-updates in the difference
        view (\(+\delta\) at \(k\), \(-\delta\) at \(k{+}1\)) \emph{and}
        still needs the shadow array.  Mixing problem types is this
        section's signature WA.
  \item \textbf{Guard or size-up:} pick one convention for \(b=n\) and keep
        it in the template (\code{if (b+1<=n)} above; or \(MX=n+2\) and no
        guard).
  \item \textbf{\(\pm u\) stress:} generate random \([a,b]\) adds and compare
        \code{query(k)} against a brute-force array for \(n\le50\), \(q\le100\)
        -- five lines of Python; catches an inverted sign at \(b{+}1\)
        instantly.
\end{itemize}
\end{hardeningbox}

% ---------------------------------------------------------------
\subsection{Common Pitfalls}
% ---------------------------------------------------------------

\begin{warningbox}[Trap 1: storing the raw difference array \emph{as if} it were the BIT]
\(D\) and the BIT over \(D\) are different objects: the second stores
\emph{block sums} of the first.  Skipping the push phase of
\(\coderef{c14-diff-build}\) (leaving \(t=D\)) makes \code{query} walk
individual \(D\) entries and return \(A_k\) only for the indices the walk
happens to touch -- e.g.\ on the sample, \(\text{query}(4)\) would answer
\(1\) (the raw \(D_4\)) instead of \(5\).  Wrong answers that ``look like
off-by-one'' are usually this.
\end{warningbox}

\begin{warningbox}[Trap 2: forgetting \(-u\) at \(b+1\)]
A one-sided difference update turns every later read into an accumulated
flood of \(+u\).  Visible only when a query position exceeds \(b\) -- many
hand samples never query there.  Say the two words out loud when writing:
\emph{enter, leave}.
\end{warningbox}

\begin{warningbox}[Trap 3: memorising the two-BIT formula with the wrong weight]
The pair is \(\big(D_j,\;(j{-}1)D_j\big)\) with
\(\mathit{prefix}(x)=x\,Q_1(x)-Q_2(x)\), \emph{or} \(\big(D_j,\;jD_j\big)\)
with \((x{+}1)Q_1-Q_2\) -- one convention, consistently.  The bug appears
when code and derivation come from different memories.  Re-derive in one
line on the page before submitting (the derivation is in
\thmref{c14-twobit}'s proof); do not recall.
\end{warningbox}

\begin{warningbox}[Trap 4: compression on a \emph{sum} that needs real values]
Compression preserves \emph{order} (counts, k-th, \(\le k\) predicates).  It
destroys \emph{magnitudes}: ``sum of values \(\le k\)'' cannot be answered
by counting compressed ranks.  Keep sums on a value-carrying structure or
store \(v\)-weighted trees -- \(\coderef{c14-twobit}\)'s second tree is the
template for weighting.
\end{warningbox}

% ---------------------------------------------------------------
\subsection{Practice Problems}
% ---------------------------------------------------------------

\begin{stuckbox}
\begin{itemize}[nosep,leftmargin=1.4em]
  \item Which side is cheap in your encoding -- update or read?  Choose the
        stored object so the \emph{frequent} op is the cheap one.
  \item Did you build the difference tree with the push phase, not raw \(D\)?
  \item Compression: are you using only \emph{order} properties?
  \item Two-BIT: which weight convention is your code actually using?
\end{itemize}
\end{stuckbox}

\begin{practiceset}
  \practice{Little Girl and Maximum Sum}{CF 276C}
    {https://codeforces.com/problemset/problem/276/C}
    {CF 1500 L3}
    {\(q\) fixed ranges, reorder the array so that the \emph{sum over the
     ranges} is maximal.  Count how many ranges cover each position --
     that is the difference-array idea of \defref{c14-diff} run once
     offline (two \(O(1)\) writes per range, one prefix sweep) -- then sort
     and match greedily.  No BIT survives the queries here (the array never
     changes), which makes the difference \emph{technique} visible without
     the tree.}
  \practice{Iahub and Xors}{CF 341D}
    {https://codeforces.com/problemset/problem/341/D}
    {CF 1900 L5}
    {A \(n\times n\) matrix, two ops: xor \(x\) into every cell of a
     subrectangle, and xor of a subrectangle.  \(\oplus\) is its own inverse
     (invertible accumulator -- the BIT licence holds), so the rectangle is
     a 2D difference encoding with four corner events per update, and a
     range-xor needs the same paired-tree trick as \thmref{c14-twobit} in
     two axes.  Budget a page of index algebra before coding; verify on a
     \(2\times2\) hand case first.}
\end{practiceset}

\paragraph{Back-cites and forwards.}
\secref{c14-core} owns the walks this section reuses without restating.
Chapter~5 (\propref{c24-prefix}) owns telescoping in its static form.
Forward: the frequency reading of compressed indices is
\secref{c14-order}; difference-on-Euler arrays is \secref{c14-trees}
(CF~383C); 2D sweep encodings return in \secref{c14-offline}.  Arbitrary
aggregates under arbitrary writes are the segment tree's -- \chref{c13},
\chref{c15}; the boundary table is above, the code is there.

% ---------------------------------------------------------------
\begin{recallbox}[Section 18.2 in 10 seconds]
Store \(D=A_{\mathit{diff}}\): interval add \(=\) two point updates
(\defref{c14-diff}); point read \(=\) prefix walk (CSES~1651,
\(\coderef{c14-1651-soln}\)).  Need range sums too?  Pair a second tree of
\((j{-}1)D_j\): \(\mathit{prefix}(x)=xQ_1-Q_2\) (\thmref{c14-twobit}) --
pick ONE weight convention.  Huge index domains: collect--sort--unique--
lower\_bound (\propref{c14-compress}); order survives, magnitudes don't.
\(\min\)-BIT legal only while writes are monotone (\propref{c14-prefixmin});
otherwise \chref{c13}.  Never confuse raw \(D\) with the BIT over \(D\).
\selftest{
  \item Which two slots does ``\(+5\) on \([3,7]\)'' touch in the difference
        tree, and what are the increments?
  \item Re-derive \(\mathit{prefix}(x)=xQ_1-Q_2\) in two lines from
        \(\sum_{j\le x}D_j(x-j+1)\).
  \item Give the workload where a lazy segment tree is the \emph{smaller}
        program than the two-BIT pair.
  \item Why does compression answer ``how many \(\le k\)'' but not ``sum of
        values \(\le k\)''?
  \item State the monotonicity hypothesis of the prefix-min BIT and the
        failure mode when it breaks.
}
\end{recallbox}
